\begin{summarybox}
  \textbf{\textcolor{CSummary}{快记｜一句口诀}}

  \[
    \boxed{\text{左移 }a:\;f(x+a)}
    \qquad
    \boxed{\text{右移 }a:\;f(x-a)}
    \quad (a>0).
  \]

  \textbf{\textcolor{CSummary}{快懂｜为什么符号相反？}}

  点向左移 $2$ 格，横坐标从 $0$ 变成 $-2$；新函数在 $-2$ 处要找回原来 $0$ 处的高度，必须用 $-2+2=0$。
  \textbf{点向左，坐标减；式子找原位置，$x$ 加。}

  \medskip
  \textbf{\textcolor{CSummary}{快用｜考场三步}}
  \begin{itemize}[leftmargin=*,itemsep=0.3em]
    \item \textbf{先判断：}左移还是右移？移几格？
    \item \textbf{再替换：}把原式\textbf{每个 $x$} 换成 $x+a$ 或 $x-a$。
    \item \textbf{后检验：}看顶点、对称轴或一个对应点是否移对；有定义域时也要跟着移动。
  \end{itemize}

  \textbf{\textcolor{CSummary}{最容易错的地方}}

  不是在整个函数后面加减常数；也不是平移只能在原定义域里进行。
  记住：\textbf{左右平移改变横坐标，不改变同一对应点的高度。}
\end{summarybox}
